% !TeX root = ../../Book3.tex
% 纲要标题：### 10.2 离散赋值环
% 纲要位置：工作记录/计划纲要.md，第 2237 行。
% 纲要细目（写作范围）：
% * 一维 Noether 正规局部整环的刻画
% * 极大理想的统一化元
% * 理想与阶函数：由商模长度恢复整数阶，并用尖点说明长度与外部赋值的相等不能单独刻画 DVR

\section{离散赋值环}
\label{b3:sec:1002}

离散赋值环最有用之处在于，所有非零理想都由同一个元素的幂生成。因此它既能按整数测量阶数，也能从局部环的有限性和整闭性识别。


\subsection{一维 Noether 正规局部整环的刻画}
\label{b3:subsec:1002:01}

\begin{theorem}\label{b3:thm:dvr-characterizations}
对非域整环 \(R\)，下列条件等价：
\begin{equivalences}
\item\label{b3:item:dvr-valuation} \(R\) 为离散赋值环；
\item\label{b3:item:dvr-principal-maximal} \(R\) 为 Noether 局部环，且其极大理想由一个元素生成；
\item\label{b3:item:dvr-normal-dimension-one} \(R\) 为一维 Noether 正规局部整环。
\end{equivalences}
它们也等价于 \(R\) 是非域的局部主理想整环。
\end{theorem}
\begin{proof}
先设 \(R=V_v\)、\(v(K^*)=\mathbb Z\)，选 \(\pi\) 满足 \(v(\pi)=1\)。
每个 \(0\ne a\in R\) 唯一写成 \(a=\pi^nu\)，其中 \(n\ge0\)、\(u\) 为单位。
非零理想中非零元素的赋值有最小值 \(n\)，于是该理想恰为 \((\pi^n)\)。所以 \(R\) 是主理想整环，且只有 \((0)\)、\((\pi)\) 两个素理想。由\cref{b3:prop:valuation-local-normal}，赋值环整闭，这同时得到\itemref{b3:item:dvr-principal-maximal}和\itemref{b3:item:dvr-normal-dimension-one}。

设\itemref{b3:item:dvr-principal-maximal}成立，\(\mathfrak m=(\pi)\)。要把非零元素写成统一化元的有限次幂乘以单位，须排除它被 \(\pi\) 的所有幂整除的情形，因此令 \(J=\bigcap_{n\ge0}\pi^nR\)。
若 \(x\in J\)，写 \(x=\pi y\)；由整环中可约去 \(\pi\)，从
\(x\in\pi^{n+1}R\) 得 \(y\in\pi^nR\)，故 \(y\in J\)。于是 \(J=\pi J\)。
Noether 性使理想 \(J\) 有限生成，且 \((\pi)\subseteq J(R)=\mathfrak m\)，所以 Nakayama 引理给出 \(J=0\)。对每个非零 \(a\in R\)，存在 \(N\) 使 \(a\notin\pi^NR\)，而 \(a\in\pi^0R\)。由于这些理想依次缩小，满足 \(a\in\pi^nR\) 的非负整数构成有限初段，故有最大值 \(n(a)\)。于是 \(a=\pi^{n(a)}u\)，且 \(u\notin\mathfrak m\)，否则 \(a\in\pi^{n(a)+1}R\)；所以 \(u\) 为单位。
两个这样的分解相乘，单位部分的乘积仍是单位，故 \(n(ab)=n(a)+n(b)\)。对 \(a,b\in R\setminus\{0\}\) 定义 \(v(a/b)=n(a)-n(b)\)。若 \(a/b=c/d\)，则 \(ad=bc\)，从而 \(n(a)+n(d)=n(b)+n(c)\)，这说明分式表示无关。乘法指数相加；对环中元素相加，提出较小的 \(\pi\) 次幂后剩余部分仍在 \(R\)，故阶数不小于两个输入阶数的最小值。对任意分式通分，便得到同样的赋值不等式。
其值群为 \(\mathbb Z\)，非负值元素恰属于 \(R\)，所以得到\itemref{b3:item:dvr-valuation}。

最后设\itemref{b3:item:dvr-normal-dimension-one}成立，记极大理想为 \(\mathfrak m\)，分式域为 \(K\)。为得到 \(\mathfrak m\) 的一个生成元，可以寻找 \(x\in K\setminus R\) 使 \(x\mathfrak m\subseteq R\)，再用整闭性证明这个理想必须含有单位。
取 \(0\ne a\in\mathfrak m\)。包含 \(a\) 的素理想非零，由维数一和局部性只能是 \(\mathfrak m\)，所以 \(\sqrt{aR}=\mathfrak m\)，故某个
\(\mathfrak m^n\subseteq aR\)。这里用到 \(\mathfrak m\) 有限生成：每个生成元的某个幂属于 \(aR\)，展开充分高次的乘积即可。
取最小的 \(n\ge1\)，再取 \(b\in\mathfrak m^{n-1}\setminus aR\)。这样的 \(b\) 存在：当 \(n=1\) 时，\(R\not\subseteq aR\)；当 \(n>1\) 时，这是 \(n\) 的最小性。
于是 \(x=b/a\notin R\)，而 \(b\mathfrak m\subseteq\mathfrak m^n\subseteq aR\)，正好给出 \(x\mathfrak m\subseteq R\)。
若 \(x\mathfrak m\subseteq\mathfrak m\)，乘以 \(x\) 就定义了有限生成 \(R\) 模 \(\mathfrak m\) 上的 \(R\) 线性自同态 \(u_x:y\mapsto xy\)。对它使用\cref{b3:thm:determinant-trick}，取该定理中的理想为 \(R\)，得到首一多项式 \(P(T)\in R[T]\) 使 \(P(u_x)=0\)。因此对每个 \(y\in\mathfrak m\) 都有 \(P(x)y=0\)。这里的忠实性来自整环中 \(\mathfrak m\ne0\)：取 \(0\ne y\in\mathfrak m\)，在 \(K\) 中约去 \(y\)，得到 \(P(x)=0\)，故 \(x\) 在 \(R\) 上整。
正规性将推出 \(x\in R\)，矛盾。
所以理想 \(x\mathfrak m\) 含有单位，从而 \(x\mathfrak m=R\)。选 \(\pi\in\mathfrak m\) 满足 \(x\pi=1\)。
每个 \(y\in\mathfrak m\) 都满足 \(xy\in R\)，于是 \(y=\pi(xy)\)，得到 \(\mathfrak m=(\pi)\)，即\itemref{b3:item:dvr-principal-maximal}。
局部主理想整环满足\itemref{b3:item:dvr-principal-maximal}，而第一段已证明\itemref{b3:item:dvr-valuation}使每个理想主生成，故最后一个等价条件也成立。
\end{proof}

Noether 性使证明中的 \(J\) 和 \(\mathfrak m\) 都有限生成：前者用于 Nakayama 引理，后者用于从 \(\sqrt{aR}=\mathfrak m\) 得到幂包含关系，并应用行列式技巧。若直接删除有限性，上一节的一维非离散赋值环就给出反例。

\subsection{极大理想的统一化元}
\label{b3:subsec:1002:02}

\begin{definition}\label{b3:def:uniformizer}
设 \((R,\mathfrak m)\) 为 DVR，\(K=\operatorname{Frac}(R)\)。若元素 \(\pi\in R\) 满足 \(\mathfrak m=(\pi)\)，则称其为
\term[tongyihuayuan]{统一化元}[uniformizer]。将相应赋值的赋值群与通常排序的 \(\mathbb Z\) 认同，并令 \(v(\pi)=1\)，所得满射 \(v:K^*\to\mathbb Z\) 称为 \(R\) 的\term[guiyihua-lisan-fuzhi]{归一化离散赋值}[normalized discrete valuation]。
\end{definition}
\symidx{\(\pi\)：离散赋值环的统一化元}

统一化元并不唯一：\(\pi'\) 是另一个统一化元，当且仅当 \(\pi'=u\pi\)、\(u\in R^*\)。
所以元素的展开 \(a=u\pi^n\) 中单位部分依赖选择，整数 \(n\) 则与选择无关。
对 \(\mathbb Z_{(p)}\) 可取 \(p\)，对 \(k[t]_{(t-c)}\) 可取 \(t-c\)。
归一化离散赋值确定了理想的阶数。环与形式幂级数环的关系，还涉及完备性、特征及剩余域中系数的提升。

\begin{proposition}\label{b3:prop:noether-valuation-dvr}
非域赋值环为 Noether 环，当且仅当它为 DVR。
\end{proposition}
\begin{proof}
DVR 已证明为主理想整环。反向，Noether 赋值环的极大理想有限生成，因而由\cref{b3:prop:valuation-ideal-comparison}主生成，再应用\cref{b3:thm:dvr-characterizations}。
\end{proof}

极大理想主生成本身也不能推出 Noether 性。设 \(k\) 为域，取 \(K=k(\!(u)\!)(\!(t)\!)\)，对非零
\(f=t^n\cdot a_n(u)+t^{n+1}\cdot a_{n+1}(u)+\cdots\) 定义
\(v(f)=(n,\operatorname{ord}_u a_n)\)，赋值群按字典序排序的 \(\mathbb Z^2\)。
首项乘法和加法直接验证赋值公理。最小正值是 \((0,1)\)，所以极大理想由 \(u\) 生成；但零元素及赋值第一坐标为正的元素组成非零素理想
\(\mathfrak p\subsetneq\mathfrak m\)。对任意 \(0\ne a\in\mathfrak p\)，都能取非负整数 \(r\) 使 \(v(t/u^r)=(1,-r)\prec v(a)\)，所以 \(\mathfrak p\) 中没有赋值最小的非零元素。
若 \(\mathfrak p\) 有限生成，\cref{b3:prop:valuation-ideal-comparison}会使它主生成，其生成元的赋值就是非零元素的最小赋值，矛盾。因此 \(\mathfrak p\) 不是有限生成理想，这个环不是 Noether 环，也不是 DVR。这正是上一节字典序 \(\mathbb Z^2\) 对应图中的中间素理想。

\subsection{理想与阶函数}
\label{b3:subsec:1002:03}

\begin{proposition}\label{b3:prop:dvr-order-ideals}
在离散赋值环 \((R,\mathfrak m)\) 中，每个非零理想唯一写成 \(\mathfrak m^n\)，其中 \(n\ge0\) 为整数。
对非负整数 \(a,b\)，有
\[
\mathfrak m^a\mathfrak m^b=\mathfrak m^{a+b},\qquad
\mathfrak m^a+\mathfrak m^b=\mathfrak m^{\min(a,b)},\qquad
\mathfrak m^a\cap\mathfrak m^b=\mathfrak m^{\max(a,b)}.
\]
并且 \(\ell_R(R/\mathfrak m^n)=n\)。若 \(v\) 为 \(R\) 的归一化离散赋值，则对每个 \(0\ne a\in R\) 有
\[
v(a)=\ell_R(R/aR).
\]
不过，这一长度等式对某些元素成立并不能反过来刻画 DVR。对任意域 \(k\) 及不定元 \(t\)，尖点局部环 \(C=k[t^2,t^3]_{(t^2,t^3)}\) 满足
\[
\ell_C(C/t^2C)=2=\operatorname{ord}_0(t^2),\qquad
\ell_C(C/t^3C)=3=\operatorname{ord}_0(t^3),
\]
但 \(t\) 在 \(C\) 上整而不属于 \(C\)，所以 \(C\) 不是 DVR。这里 \(\operatorname{ord}_0\) 是 \(k(t)\) 上的零点阶赋值。
\end{proposition}
\begin{proof}
取统一化元 \(\pi\)。理想形式已在刻画定理中证明；\(\pi^n\notin\mathfrak m^{n+1}\) 保证指数唯一。
乘积公式由生成元相乘得到。指数大小与理想包含方向相反：\(a\le b\) 时 \(\mathfrak m^a\supseteq\mathfrak m^b\)，所以和取较大的理想，对应较小的指数；交取较小的理想，对应较大的指数。
对每个 \(i\ge0\)，乘以 \(\pi^i\) 给出
\(R/\mathfrak m\simeq\mathfrak m^i/\mathfrak m^{i+1}\)：满射显然，核为 \(\mathfrak m\)，因为可在整环中约去 \(\pi^i\)。
因此 \(R/\mathfrak m^n\) 的幂滤过有恰好 \(n\) 个简单商，长度为 \(n\)。
对 \(0\ne a\in R\)，写 \(a=u\pi^n\)，其中 \(u\) 为单位。于是 \(aR=\mathfrak m^n\)，从而 \(v(a)=n=\ell_R(R/\mathfrak m^n)=\ell_R(R/aR)\)。

对尖点，记 \(A=k[t^2,t^3]\)。其中出现的幂指数恰为 \(0,2,3,4,\ldots\)，所以 \(A/t^2A\) 的 \(k\) 基为 \(1,t^3\)，\(A/t^3A\) 的 \(k\) 基为 \(1,t^2,t^4\)。这两个商环分别同构于 \(k[z]/(z^2)\) 和 \(k[z]/(z^3)\)，其中 \(z\) 分别对应 \(t^3\) 和 \(t^2\)。常数项非零的元素在两商环中已经可逆，因此局部化不改变它们。沿 \(z\) 的幂所得的链分别有二个和三个简单商，均为剩余域 \(k\)，故作为 \(C\) 模的长度分别为二和三。
另一方面，\(t=t^3/t^2\) 属于 \(\operatorname{Frac}(C)=k(t)\)，并满足首一方程 \(T^2-t^2=0\)，故在 \(C\) 上整。若 \(t\in C\)，可写 \(t=f/g\)，其中 \(f,g\in A\)、\(g(0)\ne0\)。但 \(tg=f\) 的左端一次项系数为 \(g(0)\ne0\)，右端没有一次项，矛盾。因此 \(C\) 不整闭，不能为 DVR。
\end{proof}

在 DVR 上，整数阶因此可以完全由商模的合成层数读取。在一维正规曲线中，这将给出零点重数的一个局部代数描述。

\ProseParagraphBreak
一般赋值环的理想已按包含关系排成全序；Noether 条件则使非域赋值环成为 DVR。下一节的 Dedekind 整环把一维正规条件放在整个环上，在各个非零素理想处得到 DVR，进而合成全局的理想分解。\cref{b3:tab:valuation-dvr-dedekind}并列这些条件与结构。

\begin{center}
\begin{minipage}{\textwidth}
% 此表随正文排放，避免浮到后续定理与证明之间；沿用标准表题及编号。
\makeatletter
\def\@captype{table}
\makeatother
\centering
\caption[赋值环、DVR 与 Dedekind 整环]{赋值环、DVR 与 Dedekind 整环。赋值环行中，\(v:K^*\to\Gamma\) 为域 \(K\) 上的满赋值，\(V_v\) 为其赋值环。Dedekind 整环行限定为非域，其定义及全局理想分解见\secref{b3:sec:1003}。}
\label{b3:tab:valuation-dvr-dedekind}
\begin{tabularx}{\textwidth}{@{}>{\raggedright\arraybackslash}p{0.18\textwidth}>{\raggedright\arraybackslash}p{0.23\textwidth}>{\raggedright\arraybackslash}p{0.14\textwidth}>{\raggedright\arraybackslash}X@{}}
\toprule
对象 & 局部性与有限性 & Krull 维数 & 整除或理想结构\\
\midrule
一般赋值环 \(V_v\) & 局部整闭整环，未必 Noether & \(\operatorname{rank}(v)\) & 理想按包含关系全序；每个有限生成理想主生成\\
秩一赋值环 & 局部整闭整环，未必 Noether & \(1\) & 赋值群未必同构于 \(\mathbb Z\)，极大理想未必主生成\\
DVR & Noether 局部整闭整环，非域 & \(1\) & 值群同构于 \(\mathbb Z\)，非零理想为极大理想的幂；等价于一维 Noether 正规局部整环\\
非域 Dedekind 整环 & Noether 整闭整环，未必局部 & \(1\) & 非零素理想处的局部环为 DVR；非零理想有唯一素理想分解\\
\bottomrule
\end{tabularx}
\end{minipage}
\end{center}
